\hypertarget{appendix-e2d7b4f2161d}{}\label{referencestart:appendix-e2d7b4f2161d}
\section{Finite differences and step-size balance}\label{appendix-e2d7b4f2161d-finite-differences-and-step-size-balance}

\hypertarget{appendix-e2d7b4f2161d-finite-differences-and-step-size-balance}{}

\textbf{Curriculum depth:} Developed · \textbf{First seen:} Phase 9 · \textbf{Latest development:} Phase 9

\begin{allmightroledefinition}

\subsection{Definition and mechanism}\label{appendix-e2d7b4f2161d-definition-and-mechanism}

\hypertarget{appendix-e2d7b4f2161d-definition-and-mechanism}{}

Finite differences estimate derivatives from sampled function values. Forward and backward differences use one-sided samples and have first-order truncation error; central differences use samples on both sides and usually have second-order truncation error for smooth functions. Making the step arbitrarily small is not always better: subtractive cancellation and evaluation noise grow as the step shrinks, while truncation error grows when it is too large. The useful step balances these effects and depends on scale, precision, and function noise. Finite differences are valuable for black-box sensitivities and gradient checks, but they are not automatic differentiation and can fail near discontinuities or branch changes. A numerical gradient check should compare directional derivatives, use scale-aware steps, and avoid interpreting one tolerance pass as proof that every derivative path is correct.

\end{allmightroledefinition}

\subsection{Prerequisites}\label{appendix-e2d7b4f2161d-prerequisites}

\hypertarget{appendix-e2d7b4f2161d-prerequisites}{}

\begin{itemize}
\tightlist
\item
  \hyperref[appendix-7bc181935c42]{Derivative as a local linear map}
\item
  \href{../01-Bits-Programs-and-Reliable-Systems/main.pdf\#nameddest=appendix-3ae74e89b4e6}{Absolute and relative error (Bits, Programs, and Reliable Systems, Reference notes: Absolute and relative error)}
\end{itemize}

\subsection{Used by}\label{appendix-e2d7b4f2161d-used-by}

\hypertarget{appendix-e2d7b4f2161d-used-by}{}

\begin{itemize}
\item
  This node is a current endpoint in the supplied Phase 9--10 path.
\item
  \hyperref[appendix-aca80b6bf17a]{Convergence order and Richardson extrapolation} \#\# Full teaching explanation
\item
  \href{../02-Mathematics-for-Visual-Computing/main.pdf\#nameddest=02-foundations-ii-ch022-numerical-differentiation-and-automatic-differentiation}{Phase 9 teaching section (Mathematics for Visual Computing, Ch. 6)}
\item
  \href{../01-Bits-Programs-and-Reliable-Systems/main.pdf\#nameddest=appendix-adf48f5dc5c1}{Visual Recall/09 - Calculus and Optimization - Visual Recall.md (Bits, Programs, and Reliable Systems, Reference notes: 09 - Calculus and Optimization - Visual Recall)} · visual recall
\end{itemize}

\subsection{Source coverage}\label{appendix-e2d7b4f2161d-source-coverage}

\hypertarget{appendix-e2d7b4f2161d-source-coverage}{}

Supplied lesson parts: 847 (T03 · M052), 848 (T03 · M052), 849 (T03 · M052), 850 (T03 · M052), 851 (T03 · M052), 884 (T02 · M054).
